\chapter*{GLOSARIO Y ACRÓNIMOS}
\addcontentsline{toc}{chapter}{Glosario y Acrónimos}
\markboth{GLOSARIO Y ACRÓNIMOS}{GLOSARIO Y ACRÓNIMOS}

\begin{description}
    \item[API] Application Programming Interface (Interfaz de Programación de Aplicaciones). Conjunto de definiciones y protocolos que se utiliza para desarrollar e integrar el software de las aplicaciones.
    \item[APK] Android Package Kit. Formato de archivo de paquete utilizado por el sistema operativo Android para la distribución e instalación de aplicaciones móviles.
    \item[DAO] Data Access Object (Objeto de Acceso a Datos). Patrón de diseño que proporciona una interfaz abstracta para algún tipo de base de datos u otro mecanismo de persistencia.
    \item[HTTP] Hypertext Transfer Protocol (Protocolo de Transferencia de Hipertexto). Protocolo de comunicación que permite las transferencias de información en la World Wide Web.
    \item[IDE] Integrated Development Environment (Entorno de Desarrollo Integrado). Aplicación informática que proporciona servicios integrales para facilitarle al desarrollador o programador el desarrollo de software.
    \item[JSON] JavaScript Object Notation (Notación de Objetos de JavaScript). Formato ligero de intercambio de datos, fácil de leer y escribir para humanos y fácil de analizar y generar para máquinas.
    \item[JVM] Java Virtual Machine (Máquina Virtual de Java). Máquina virtual de proceso nativo, es decir, ejecutable en una plataforma específica, capaz de interpretar y ejecutar instrucciones expresadas en un código binario especial (el bytecode de Java).
    \item[LLM] Large Language Model (Gran Modelo de Lenguaje). Modelo de inteligencia artificial entrenado en grandes cantidades de datos de texto para entender y generar lenguaje humano.
    \item[MVVM] Model-View-ViewModel (Modelo-Vista-Modelo de Vista). Patrón de arquitectura de software que facilita la separación del desarrollo de la interfaz gráfica de usuario de la lógica de negocio o back-end.
    \item[PDF] Portable Document Format (Formato de Documento Portátil). Formato de almacenamiento para documentos digitales independiente de plataformas de software o hardware.
    \item[PII] Personally Identifiable Information (Información de Identificación Personal). Cualquier información que pueda utilizarse para distinguir o rastrear la identidad de una persona.
    \item[SDK] Software Development Kit (Kit de Desarrollo de Software). Conjunto de herramientas de desarrollo de software que permite a los desarrolladores crear aplicaciones para un sistema concreto.
    \item[SQL] Structured Query Language (Lenguaje de Consulta Estructurado). Lenguaje de dominio específico utilizado en programación, diseñado para administrar, y recuperar información de sistemas de gestión de bases de datos relacionales.
    \item[UI] User Interface (Interfaz de Usuario). Espacio donde se producen las interacciones entre seres humanos y máquinas.
    \item[URL] Uniform Resource Locator (Localizador Uniforme de Recursos). Dirección única que se utiliza para localizar un recurso en Internet.
    \item[UX] User Experience (Experiencia de Usuario). Conjunto de factores y elementos relativos a la interacción del usuario con un entorno o dispositivo concretos, cuyo resultado es la generación de una percepción positiva o negativa de dicho servicio, producto o dispositivo.
    \item[XML] eXtensible Markup Language (Lenguaje de Marcado Extensible). Lenguaje de marcado que define un conjunto de reglas para la codificación de documentos.
\end{description}
